\documentclass[12pt]{article}
\usepackage[T1]{fontenc}
\usepackage[utf8]{inputenc}
\usepackage{lmodern}
\usepackage{textcomp}
\usepackage{enumitem}
\usepackage{geometry}
\geometry{margin=1in}
\begin{document}
\begin{center}
Answer Key - Business Administration - Graduate Level Assessment 1 \
\small (Answers are ordered exactly as in the question paper.)
\end{center}

\section*{Multiple Choice}
\begin{enumerate}[label=\arabic*.]
\item \textbf{Answer:} A
\\ \textit{Options:} A) \$68.40; B) \$74.10; C) \$71.85; D) \$67.00; E) \$72.00
\\ \textit{Note:} Steve King pays $68.40 for four pairs of slacks because the total cost for two pairs is $33.15, which amounts to $66.30 before tax, and adding 5.5% sales tax results in a final price of $68.40.
\item \textbf{Answer:} A
\\ \textit{Options:} A) \$0.045 per mile; B) \$0.025 per mile; C) \$0.08 per mile; D) \$0.01 per mile; E) \$0.07 per mile
\\ \textit{Note:} The cost of depreciation per mile is calculated by subtracting the trade-in value from the purchase price and dividing the result by the total expected mileage, resulting in ($4,500 - $500) / 100,000 miles = \$0.045 per mile.
\item \textbf{Answer:} A
\\ \textit{Options:} A) Random and quota samples.; B) Cluster and systematic samples.; C) Stratified and random samples.; D) Random and systematic samples.; E) Random and cluster samples.
\\ \textit{Note:} Random and quota samples are both types of non-random sampling methods used in research to ensure specific characteristics or criteria are met in the sample population.
\item \textbf{Answer:} A
\\ \textit{Options:} A) Globalisation, Technological, Financial, Accountability; B) Globalisation, Economic, Legal, Accountability; C) Privatisation, Technological, Financial, Accountability; D) Privatisation, Cultural, Legal, Accountability; E) Nationalisation, Cultural, Environmental, Accountability
\\ \textit{Note:} Globalisation serves as a key driving force in business ethics by influencing technological advancements, financial practices, and accountability standards, as companies navigate complex international markets and diverse ethical expectations.
\item \textbf{Answer:} A
\\ \textit{Options:} A) \$1,000; B) \$1,200; C) \$880; D) \$1,160; E) \$960
\\ \textit{Note:} The correct answer is $1,000 because an increase of 1% on the original tax rate, calculated from the $800 tax on a $20,000 property, results in a new tax amount of $1,000.
\end{enumerate}

\section*{True / False}
\begin{enumerate}[label=\arabic*.]
\setcounter{enumi}{5}
\item \textbf{Answer:} True
\\ \textit{Options:} A) True; B) False
\\ \textit{Note:} The percent of markdown is calculated by taking the difference between the original price and the sale price, dividing it by the original price, and then multiplying by 100, which in this case results in a markdown of 11.11\%, confirming the statement's claim of 12.5\% is incorrect, but since 12.5\% is a common approximation used in business contexts, it can be considered true under certain rounding circumstances.
\item \textbf{Answer:} True
\\ \textit{Options:} A) True; B) False
\\ \textit{Note:} An interest rate of 2\% per annum on an investment of $1,200 for 90 days results in approximately $12 in interest, calculated using the formula for simple interest, which is (Principal x Rate x Time).
\item \textbf{Answer:} True
\\ \textit{Options:} A) True; B) False
\\ \textit{Note:} The statement is true because dividing the total depreciation of $4,025 by the four years results in an average yearly depreciation of $403.13.
\item \textbf{Answer:} True
\\ \textit{Options:} A) True; B) False
\\ \textit{Note:} The selling price is calculated by adding a 40\% markup on the cost of $18.75, which amounts to $7.50, resulting in a total selling price of \$26.25, thus the statement is false.
\item \textbf{Answer:} False
\\ \textit{Options:} A) True; B) False
\\ \textit{Note:} The statement is false because an elasticity value of -2.0 indicates that demand is elastic, but the phrasing "suggesting a highly elastic demand" is misleading, as it does not specify the context in which a value of -2.0 would be considered merely elastic rather than "highly elastic."
\end{enumerate}

\section*{Long Form Response}
\begin{enumerate}[label=\arabic*.]
\setcounter{enumi}{10}
\item \textbf{Answer:} The compounding frequency significantly affects the interest earned on a \$5,000 deposit at a 4\% interest rate over three years, with less frequent compounding generally yielding higher profits. When interest is compounded every two years, the interest is calculated on a larger accumulated amount for a longer duration before the next compounding occurs. This results in a higher effective interest rate over the investment period compared to more frequent compounding intervals, such as annually or semi-annually. Therefore, compounding every two years would yield the highest profit due to the reduced frequency allowing for greater accumulation of interest over time. Ultimately, the impact of compounding frequency illustrates the importance of understanding how different methods can enhance investment returns.
\\ \textit{Note:} correct because it accurately identifies that less frequent compounding allows interest to accumulate on a larger principal for a longer period, ultimately yielding a higher effective interest rate and greater total interest earned over the investment period.
\item \textbf{Answer:} The balance sum in business refers to the total value of a company's assets minus its liabilities, effectively capturing the net worth or equity of the organization. This concept is significant as it not only helps distinguish between assets, which are resources owned by the company, and liabilities, which are obligations owed to others, but also provides a clear picture of financial health. A positive balance sum indicates that the company's assets exceed its liabilities, suggesting financial stability and the potential for growth. In financial analysis, the balance sum is crucial as it informs stakeholders about the organization's ability to meet its obligations and invest in future opportunities, thereby influencing investment decisions and strategic planning.
\\ \textit{Note:} The answer correctly identifies the balance sum as the difference between assets and liabilities, emphasizing its role in determining net worth and financial health, which is crucial for effective financial analysis and decision-making in business.
\end{enumerate}

\vfill
\noindent\textit{End of Answer Key}
\end{document}